\documentclass[11pt,a4paper]{article}
\usepackage[margin=2.2cm]{geometry}
\usepackage{amsmath,booktabs,longtable,array,hyperref}
\setlength{\parskip}{0.5em}\setlength{\parindent}{0pt}
\title{Live Caution Sensor and End-to-End Episode Walkthrough\\\large Camera-pipeline training of OV-Split Camera and Live Caution Camera}
\date{14 November 2026 --- working document, nothing in it is committed or pushed}
\author{}
\begin{document}\maketitle
\tableofcontents

\section{Live caution sensor}
\subsection{What it does}
Each simulation step the sensor looks at the front-left camera picture \emph{after} the synthetic weather filter (the same picture the stereo and YOLO modules receive) and outputs a caution level $c\in[0,1]$. It reads pixels only: no weather label, no simulator state. The driving network is unchanged; it receives this $c$ in the input slot that previously received the label-derived value.

\subsection{Formula}
The picture is shrunk to quarter size, converted to gray $g$ (mean of RGB, range 0--1). Seven numbers are computed:
\begin{center}\small\begin{tabular}{@{}llp{8.2cm}@{}}\toprule
$k$ & feature & definition\\\midrule
1 & log\_lapvar & $\log(\mathrm{Var}[\nabla^2 g]+10^{-8})$ (sharpness; Laplacian)\\
2 & grad & mean Sobel gradient magnitude $\overline{\sqrt{g_x^2+g_y^2}}$\\
3 & contrast & standard deviation of $g$\\
4 & dark & mean of the dark channel (per-pixel minimum over RGB, then $7{\times}7$ erosion)\\
5 & bright & mean of $g$\\
6 & sat & mean of $(\max_{RGB}-\min_{RGB})/(\max_{RGB}+10^{-3})$\\
7 & edge & share of pixels with gradient magnitude $>0.15$\\\bottomrule\end{tabular}\end{center}
With $\mathbf{f}$ the feature vector, $\mu,\sigma$ the per-feature mean and standard deviation of the fitting frames, and weights $\mathbf{w}$, bias $b$:
\[
\tilde c_t=\mathrm{clip}\!\Big(\mathbf{w}^{\top}\tfrac{\mathbf f_t-\mu}{\sigma}+b,\;0,\;1\Big),\qquad
\bar c_t=(1-\alpha)\,\bar c_{t-1}+\alpha\,\tilde c_t,\ \ \bar c_0=\tilde c_0,\ \ \alpha=0.15,
\]
\[
\hat c_t=\begin{cases}0&\bar c_t<0.1997\ \text{(dead-zone)}\\ \bar c_t&\text{otherwise.}\end{cases}
\]
The smoother is reset at the start of every episode. $\hat c_t$ is what the policy sees.

\subsection{Fitted constants}
$\mathbf w=(-0.577,\,-0.420,\,0.173,\,1.397,\,-1.155,\,0.259,\,1.189)$, $b=0.3986$, $\mu=(-5.915,\,0.1258,\,0.1597,\,0.3019,\,0.386,\,0.2372,\,0.1815)$, $\sigma=(1.441,\,0.0758,\,0.0285,\,0.1545,\,0.1488,\,0.0802,\,0.1209)$. Stored in \texttt{bev\_rl/adaptive/results/visibility\_sensor/weights.json}; code in \texttt{bev\_rl/adaptive/visibility\_sensor.py}.

\subsection{Calibration targets}
Ridge regression (penalty 1.0 on standardised features) of the feature vector onto the target
\[ c^{\star}=\mathrm{clip}\big(1-\mathrm{keep}(\text{weather},\text{level}),0,1\big), \]
where $\mathrm{keep}$ is the YOLO retention measured offline and stored in \texttt{retention\_table.json} (linear interpolation between measured levels). This is the same quantity the original caution model was given. The weather label is therefore used \textbf{only as the regression target in the offline fit}. Target values: clear 0, fog 0.3 $\to$ 0, fog 0.6 $\to$ 0, fog 1.0 $\to$ 0.733, rain 0.2 $\to$ 0.163, rain 0.45 $\to$ 0.491, rain 0.6 $\to$ 0.842, rain 1.0 $\to$ 0.960.

\subsection{Training (fitting) scenes}
42 MetaDrive scenes, seeds 2000--2041 (traffic density 0.1). In each, 5 frames (steps 8, 16, 24, 32, 40) were taken and each frame rendered under 8 weather settings: 42 scenes $\times$ 5 frames $\times$ 8 settings = 1\,680 frames. Weather noise seed: $20261114+100003\cdot\text{seed}+t$.

\subsection{Held-out scenes (cross-validation)}
5-fold cross-validation \emph{by scene}: the 42 scenes were split into 5 groups, and each group's frames were predicted by a fit that never saw that group's scenes. Results (raw estimate $\tilde c$, before smoothing and dead-zone):
\begin{center}\small\begin{tabular}{@{}lrrrrr@{}}\toprule
weather & target $c^\star$ & mean est. & std & 5th pct & 95th pct\\\midrule
clear & 0.000 & 0.064 & 0.050 & 0.000 & 0.177\\
fog 0.3 & 0.000 & 0.003 & 0.011 & 0.000 & 0.031\\
fog 0.6 & 0.000 & 0.196 & 0.040 & 0.149 & 0.270\\
fog 1.0 & 0.733 & 0.619 & 0.039 & 0.557 & 0.673\\
rain 0.2 & 0.163 & 0.097 & 0.040 & 0.020 & 0.163\\
rain 0.45 & 0.491 & 0.536 & 0.044 & 0.462 & 0.605\\
rain 0.6 & 0.842 & 0.788 & 0.051 & 0.706 & 0.874\\
rain 1.0 & 0.960 & 0.959 & 0.042 & 0.884 & 1.000\\\bottomrule\end{tabular}\end{center}
Overall mean absolute error 0.074, correlation 0.968. The dead-zone 0.1997 is the 99th percentile of the cross-validated clear-weather estimates.

\subsection{Were evaluation images used? --- confirmation and honest limits}
\textbf{Confirmed:} the sensor was fitted only on frames from seeds 2000--2041 (checked in the stored frame files: 42 distinct seeds, min 2000, max 2041). Validation seeds 3100--3139 and evaluation seeds 5000--5199 were never rendered for the sensor, never used for the dead-zone, smoothing constant or weights. The smoothing constant $\alpha=0.15$ was chosen by hand before any evaluation and was not tuned.

\textbf{Limits you should state if this goes in a paper:}
\begin{itemize}
\item The ``held-out'' numbers are cross-validation \emph{inside} the 42 fitting scenes, not a separate test set. A fully separate test (e.g. seeds 3100--3139) has not yet been run; it can be done cheaply before evaluation.
\item The \emph{weather settings} at evaluation (fog 1.0, rain 0.6, rain 1.0) are the same settings that appear in the fit; only the scenes differ. The sensor was not tested on an unseen weather level.
\item Fog 0.6 looks hazy but has target 0 (YOLO still detects), and the sensor reads 0.196 for it, right at the dead-zone: it is ambiguous there.
\item The sensor reads the \emph{synthetic filter's} appearance. It has never seen real fog or rain.
\item Its error is largest for fog 1.0 (0.62 vs 0.73) and rain 0.2 (0.10 vs 0.16): it slightly under-estimates caution there.
\end{itemize}

\section{What happens in one training episode (DAgger data collection)}
Setting: MetaDrive 0.4.3, one seed from 2000--2999, traffic density $d\in\{0,0.05,0.1,0.2,0.3\}$, horizon 600 steps (10 Hz = 60 s), weather drawn from the 10-item training mix (clear $\times4$, fog 0.6, fog 1.0, rain 0.2/0.45/0.6/1.0). 8 episodes per density per round, 5 rounds (lean).

\subsection{At reset}
The scenario is created from the seed; the weather label is fixed for the episode; the teacher is built; for the caution model the teacher's caution is $c=c^{\star}(\text{weather})$ (privileged, label-derived); the sensor smoother is reset.

\subsection{Every step (10 Hz)}
\begin{enumerate}
\item \textbf{Cameras.} Five RGB cameras, 1200$\times$900, 60$^\circ$ field of view, height 1.4 m, stereo baseline 0.5 m, focal length $f=\tfrac{600}{\tan 30^\circ}=1039.23$ px.
\item \textbf{Weather filter} (existing module, unchanged) on all pictures; the two front pictures are filtered together so the stereo pair stays consistent.
\item \textbf{Live caution sensor} (Section 1) on the front-left picture $\to\hat c_t$.
\item \textbf{Stereo depth} (existing, unchanged): StereoSGBM, 192 disparities, $5{\times}5$ block; depth $Z=\dfrac{f\,B}{d}$ with disparity $d$ in pixels and baseline $B=0.5$ m.
\item \textbf{YOLOv8n} (existing, unchanged) on front, left, right and rear pictures $\to$ boxes, class, confidence.
\item \textbf{Detection to ego position} (existing \texttt{objects\_from\_detections}, unchanged): box and depth $\to(x_{\text{right}},y_{\text{fwd}})$; the method (stereo or monocular) is recorded per object.
\item \textbf{BEV grid builder} (the only module modified; \texttt{camera\_grid\_builder.py}), per object:
 \begin{itemize}
 \item side cameras: $(x,y)\mapsto(-x,-y)$, then $\begin{bmatrix}x'\\y'\end{bmatrix}=\begin{bmatrix}1.20&-0.04&-1.50\\0.02&1.67&0.27\end{bmatrix}\begin{bmatrix}x\\y\\1\end{bmatrix}$;
 \item front stereo: same form with $\begin{bmatrix}1.15&0.02&-0.07\\0.11&1.03&0.87\end{bmatrix}$;
 \item rear: shift outward by half a footprint along the line of sight (not validated);
 \item every object gets the mean true footprint $2.0\times4.9$ m;
 \item duplicates merged if closer than 3.5 m forward and 1.6 m sideways (stereo and higher confidence kept).
 \end{itemize}
 The affine constants were fitted by least squares on 85 controlled placements (seeds 2000--2002; leave-one-scene-out median error front 1.58$\to$0.91 m, side 11.4$\to$1.3 m). Objects are then rasterised by the original \texttt{build\_bev\_grid} into a $64\times64$ grid (free 0, centreline 0.5, boundary 0.8, vehicle 1.0; ego at row 56, column 32) and the HD-map layers are added.
\item \textbf{Scalar state} from the simulator's ego and map: speed, progress, goal distance (the policy keeps indices 0, 1, 5, mask $[1,1,0,0,0,1]$), standardised as $(s-\mu)/\sigma$ with the original checkpoint's constants.
\item \textbf{Student policy} (architecture unchanged): CNN on the grid + MLP on scalars (+ $\hat c_t$ for the caution model), fused to 160 units, two tanh outputs (steering, throttle/brake).
\item \textbf{Teacher label} (privileged, exact simulator state): see Section 2.4. This is stored as the target $y_t$.
\item \textbf{Who drives:} the teacher with probability $\beta$ ($\beta=1$ in round 0, $0.5^{t}$ in round $t$), otherwise the student. Gaussian noise $\mathcal N(0,0.1^2)$ is added to the executed action, then it is clipped to $[-1,1]$. The stored label is the clean teacher action.
\item \textbf{Speed-up (deviation from the original recipe):} on teacher-driven odd steps the pipeline (depth, YOLO, builder) is skipped and the sample is not recorded; the sensor still updates. Student-driven steps always run the full pipeline.
\end{enumerate}
Stored per recorded step: grid (rounded to 0.1), scalars, teacher action, $\hat c_t$.

\subsection{Teachers (privileged)}
\begin{itemize}
\item \textbf{IDM base:} MetaDrive's Intelligent Driver Model, acceleration $a=a_{\max}\big[1-(v/v_0)^4-(s^{*}/s)^2\big]$, $s^{*}=d_0+vT+\tfrac{v\Delta v}{2\sqrt{ab}}$.
\item \textbf{Caution parameters:} target speed $v_0=30\,(1-c/3)$ km/h and wanted gap $d_0=10+15c$ m.
\item \textbf{SafeTeacher layer:} lane change off; every other vehicle predicted at constant velocity over 3 s; time-to-collision $t$. If $t<1.3$ s: full brake. If $1.3\le t<2.4$ s: brake $a=\min\!\big(a_{\text{IDM}},\,-0.2-0.8\,\tfrac{2.4-t}{2.4-1.3}\big)$.
\item \textbf{Curve guard:} $v_{\text{cap}}=\max\!\big(3.5,\sqrt{a_{\text{lat}}/\kappa_{\max}}\big)$ m/s with $a_{\text{lat}}=4.0$ m/s$^2$, $\kappa_{\max}$ the largest heading change per metre over the next $\max(15\text{ m},3v)$ of route; target speed $=\min(v_0,v_{\text{cap}})$.
\item OV-Split Camera: Teacher A = SafeTeacher + guard ($c=0$). Live Caution Camera: Teacher B = caution IDM + SafeTeacher + guard ($c=c^{\star}$).
\end{itemize}

\subsection{After each round}
Retrain the same network on \emph{all} data so far: loss $\frac1N\sum\|\pi_\theta(\text{grid},\text{scalars},\hat c)-y\|^2$, Adam, learning rate $10^{-3}$, batch 128; epochs 24 (round 0) / 8 later for OV-Split, 18 / 6 for caution. Initialised from the original checkpoint; originals are never overwritten. Validation (rounds 2 and 4, seeds 3100--3104, density 0.05, horizon 800): OV-Split on clear and rain 1.0; caution on clear, fog 1.0, rain 0.6, rain 1.0, with $\hat c$ from the sensor. Best round = highest (success, $-$(crash+off-road)).

\section{What happens in one evaluation episode}
Same per-step chain as training, but: held-out seeds (5000--5199; 80-episode preliminary uses 5000--5009), density 0.05, horizon 800, \emph{no teacher, no noise}: the student drives every step, and the caution model uses $\hat c_t$ from the sensor. Metrics: success / collision / off-road / timeout; completion time $=$ steps$/10$ s; estimated energy (Section 3.1); detection recall and grid agreement against the simulator grid (measurement only, never fed back).

\subsection{Energy formula (\texttt{bev\_rl/energy.py})}
Per step with speed varying linearly from $v_0$ to $v_1$ over $\Delta t$:
\[ W=\tfrac12 m(v_1^2-v_0^2)+\tfrac14 k\,\Delta t\,(v_0^3+v_0^2v_1+v_0v_1^2+v_1^3)+C_{rr}mg\tfrac{v_0+v_1}{2}\Delta t,\quad k=\tfrac12\rho C_dA,\]
$C_d=0.23$, $A=2.22$ m$^2$, $C_{rr}=0.01$, $\rho=1.225$; negative step work clipped to 0 (no regeneration). Reported as Wh/km. It is a modelled mechanical estimate, not battery electricity.

\section{Statistics used in comparisons}
Wilson 95\% intervals for rates; exact McNemar test for paired matched seeds; thresholds fixed before evaluation in the pre-registration files.

\section{Status and open items}
Training is running (OV-Split first, then Live Caution; 3 workers). Not yet done: a separate test of the sensor on seeds 3100--3139; the 80-episode evaluation; anything committed or pushed.
\end{document}
